% ================================================================
%  chapters/parts/ch09-c.tex
%  第 9 章 表面分析技术基础（分片 c）
%  源：S2（下册）§9.3–§9.5，书内 p.662—693 / PDF p.153—184
%  覆盖：9.3 UPS、9.4 AES、9.5 ISS
%  承接分片 b（9.2.5–9.2.10 XPS 数据处理与常见问题）
% ================================================================

\subsection{紫外光电子能谱（UPS）}

\begin{probout}
紫外光电子能谱（ultraviolet photoelectron spectroscopy，UPS）测量\textbf{价电子与价带}
的电子结构。它和 XPS 共用同一台谱仪的电子光学系统与分析器，唯一的差别是激发源：
用真空紫外光子（最常用 He I，$h\nu=\qty{21.22}{eV}$）代替 X 射线。
由此带来三个后果：探测深度更浅、能量分辨更高、只能看到价壳层。
因此 UPS 回答的是 XPS 回答不了的那类问题——\textbf{功函数、价带结构、吸附分子的轨道能级}。
\end{probout}

\subsubsection{原理与公式}

UPS 的基本原理与 XPS 相同，仍是 Einstein 光电方程，只是把初态从内层轨道换成价壳层。
对气相分子（原子），能量为 $h\nu$ 的光子电离出一个电子后留下可振动、转动的离子，
出射光电子的动能为

\begin{equation}\label{eq:ups-gas}
E_{\mathrm{K}}=h\nu-I-E_{\mathrm{v}}-E_{\mathrm{r}}
\end{equation}

式中 $I$ 为电离电位（把外层电子从基态激发到无穷远所需能量），$E_{\mathrm{v}}$ 为分子离子
的振动能，$E_{\mathrm{r}}$ 为转动能。$E_{\mathrm{r}}$ 至多只有几个 meV，目前无法测得；
$E_{\mathrm{v}}$ 约 \qtyrange{0.05}{0.5}{eV}，只有高分辨谱仪（分辨能力
\qtyrange{10}{25}{meV}）才能观察到振动精细结构。

对固体表面，忽略平动与转动能，价带电子的能量关系为

\begin{equation}\label{eq:ups-solid}
E_{\mathrm{K}}=h\nu-E_{\mathrm{b}}-\Phi
\end{equation}

式中 $E_{\mathrm{b}}$ 为出射光电子的结合能（相对费米能级），$\Phi$ 为逸出功。
式~\eqref{eq:ups-solid} 是 UPS 测功函数的依据。

\begin{keycond}
式~\eqref{eq:ups-gas} 要求气相、单光子电离且忽略平动转动能；式~\eqref{eq:ups-solid}
要求固体表面\textbf{清洁且导电良好}，并且样品与谱仪之间有明确的费米能级参考。
量纲检查：两式中各项均为能量，习惯以 \unit{eV} 计。$h\nu$ 由光源线型决定，
He I $\alpha$ 为 \qty{21.22}{eV}，是功函数与价带测量的常用工作点。
\end{keycond}

\begin{example}[用 UPS 测 Au 的功函数（教学设定数值）]
实验采用 He I 共振线，$h\nu=\qty{21.2}{eV}$。谱图上取费米边为 $E_{\mathrm{Fermi}}=0$，
二次电子截止边为 $E_{\mathrm{cutoff}}=\qty{16.1}{eV}$。因截止边对应动能为零的光电子，
有 $\Phi=h\nu-E_{\mathrm{cutoff}}=\qty{21.2}{eV}-\qty{16.1}{eV}=\qty{5.1}{eV}$，
与金的公认功函数一致。\pedset{本算例按 S2 §9.3.3（书内 p.669 / PDF p.160）的实测数字整理，
但请以原文为准；仪器需施加约 \qty{-5}{V} 的负偏压以使截止边落入分析器能量窗口。}
\end{example}

\subsubsection{与 XPS 的分工}

表~\ref{tab:ups-xps} 是本节最该记住的内容：UPS 与 XPS 不是替代关系，而是分工关系。

\begin{table}[htbp]
\centering
\caption{UPS 与 XPS 的分工\src{据 S2 §9.3.1--§9.3.4，书内 p.662--670 / PDF p.153--161 归纳。}}
\label{tab:ups-xps}
{\srcsize
\begin{tabularx}{\columnwidth}{@{}lXX@{}}
\toprule
对比项 & UPS & XPS \\
\midrule
激发源 & 真空紫外（He I \qty{21.22}{eV}） & X 射线（Al K$\alpha$ \qty{1486.6}{eV}） \\
本征线宽 & 约 \qty{0.01}{eV}，单色化可达 \qty{1}{meV} & 约 \qty{0.5}{eV} 量级 \\
探测对象 & 价电子、价带、表面态 & 内层电子 \\
信息深度 & 约 10 个原子层 & 约 30 个原子层 \\
价带谱质量 & 强度与分辨均优于 XPS & 可用但分辨较差 \\
化学位移 & 峰宽、不易精确测，多为峰形变化 & 内层化学位移明确，是主用途 \\
能否分辨振动结构 & 能（\qty{0.05}{eV}），单色化后甚至可分辨转动 & 不能（内层结合能远高于振转能，且源线宽过宽） \\
谱解释难度 & 较难，严谨解释需分子轨道计算 & 相对直接 \\
\bottomrule
\end{tabularx}}
\end{table}

\begin{pitfall}
\textbf{（1）「UPS 也能测化学态」是误解。}UPS 涉及成键轨道，电子属于整个分子，
谱峰宽，不易精确测化学位移，一般只能靠峰形变化判别。判断化学态仍应以 XPS 为主。
\textbf{（2）UPS 的信息深度更浅，因此对污染更敏感。}表面污染物（碳、氧）的价电子
谱线会直接叠在待测价带上。未做氩离子刻蚀清洁的 Au 价带谱，费米边与轨道能级
根本分不清。
\textbf{（3）绝缘样品特别是粉末不适合 UPS。}粉末表面存在无法去除的污染层，
且荷电无法有效中和。
\textbf{（4）谱图不直接等于态密度。}固体表面测得的光电子能量分布不直接代表状态分布，
需要通过简化模型计算才能与态密度对应。
\end{pitfall}

\begin{measure}
UPS 实验要点：分析室真空度应优于 \qty{1e-9}{mbar}；选择合适负偏压（一般约
\qty{-5}{V}）以观测截止边；气源纯度优于 $99.999\%$ 且气路短、无泄漏；
UV 束斑达 \unit{mm} 级，故样品不宜过小；导电样品与样品台须良好接地；
还需考虑 UV 对样品表面的副效应。
\src{S2 §9.3.5，书内 p.670--671 / PDF p.161--162}
\end{measure}

\subsection{俄歇电子能谱（AES）}

\begin{probout}
俄歇电子能谱（Auger electron spectroscopy，AES）用电子束激发，探测\textbf{俄歇电子}。
最关键的性质是：\textbf{俄歇电子的能量只由原子的能级结构决定，与激发源的能量和种类无关。}
这一性质使 AES 既能在固定能量处识别元素，又能用电子束聚焦获得很高的空间分辨率，
因而特别适合微区分析与薄膜、器件的失效分析。它是目前应用广度仅次于 XPS 的表面分析方法。
\end{probout}

\subsubsection{原理}

俄歇过程涉及三个电子：入射粒子（光子、电子或离子）击出内层电子形成空穴；
外层电子向空穴跃迁并释放能量；该能量以非辐射方式激发另一个电子（即俄歇电子）逃逸。
因为能量来自原子自身的能级差，所以俄歇峰的位置\textbf{不随激发源改变}——
这正是 AES 可以做成扫描微区成分分析的基础。

\begin{keycond}
俄歇电子能量在 \unit{keV} 量级，衰减长度很短，只有前几个原子层
（约 2--10 层）产生的俄歇电子能逃逸出来。因此 AES 比 XPS 更表面敏感，
但也意味着\textbf{任何表面污染都会强烈影响结果}。AES 通常只适用于导电与半导电样品，
绝缘样品必须先解决荷电问题。
\end{keycond}

\subsubsection{定量与校准}

AES 定量采用相对灵敏度因子法：

\begin{equation}\label{eq:aes-quant}
x_{\mathrm{A}}=\frac{I_{\mathrm{A}}/s_{\mathrm{A}}}{\sum_{i} I_{i}/s_{i}}
\end{equation}

式中 $x_{\mathrm{A}}$ 为元素 A 的原子分数，$I_{\mathrm{A}}$ 为 A 的 AES 强度，
$s_{\mathrm{A}}$ 为 A 的相对灵敏度因子。由于俄歇峰（直接谱）较宽，实测中一般取
\textbf{微分谱的峰对峰高度}作为峰强度。

\begin{pitfall}
\textbf{（1）灵敏度因子随激发能量而变。}电子束能量 3 keV 与 5 keV 下，
各俄歇跃迁的相对灵敏度因子分布不同，引用时必须说明激发条件，
不能跨仪器、跨束能直接比强度。
\textbf{（2）能量标尺必须校准。}早期用 Au MNN = \qty{2024}{eV}、
Ag MNN = \qty{356}{eV}、Cu LMM = \qty{920}{eV} 三点校准；
现行做法依 ISO 标准：ISO 17973（中等分辨率 AES 能量标尺校准，用于元素分析）、
ISO 17974（高分辨率 AES，用于元素与化学态）、ISO 21270（XPS 与 AES 强度标的线性）、
ISO 24236（AES 强度标的重复性与稳定性）。
\textbf{（3）重元素谱线重叠严重。}轻元素俄歇跃迁少、易识别；重元素跃迁多、易干扰。
例如 Cr 至 Cu 的 LMM 跃迁相互重叠，而这恰是钢铁材料 AES 分析的关键位置。
必须事先确认峰位不重叠，并与标准谱图对照。
\textbf{（4）绝缘样品的处理是硬约束。}可把小绝缘颗粒镶嵌在纯软金属箔（铟、铝）中、
表面加盖金属导电筛网使电子束穿孔通过、减薄样品降低绝缘性，
或用低能正离子中和表面负电荷；无论哪种，导电层都须良好接地。
\end{pitfall}

\subsection{离子散射谱（ISS）}

\begin{probout}
离子散射谱（ion scattering spectroscopy，ISS）测量\textbf{一定能量的离子在表面原子上
弹性散射后的能量}。它探测的是\textbf{最外层原子}，是本书介绍的四类表面方法中
信息深度最浅的一个，因而对「表面最外层到底由什么原子终止」这类问题最直接。
一般多功能电子能谱仪配的是低能离子散射（LEISS，离子动能 \qtyrange{200}{3000}{eV}）；
按能量还分为 MEISS（中能）与 HEIS／RBS（高能，卢瑟福背散射）。
\end{probout}

\subsubsection{原理与公式}

已知质量 $m_1$、能量 $E_0$ 的一次离子轰击表面原子（质量 $m_2$），在固定散射角 $\theta$
处测量弹性散射后的一次离子能量 $E_1$，由能量与动量守恒得

\begin{equation}\label{eq:iss-ch09}
\frac{E_1}{E_0}=\left[\frac{\cos\theta\pm\sqrt{(m_2/m_1)^2-\sin^2\theta}}{1+m_2/m_1}\right]^{2}
\end{equation}

\begin{keycond}
式~\eqref{eq:iss-ch09} 成立的前提是\textbf{单次二元弹性碰撞}（表面足够稀薄、无多重散射）
且忽略靶原子反冲的能量损失以外的过程。极限情况检查：当 $\theta=90^{\circ}$ 时
$\cos\theta=0$，取正号分支化简得
$E_1/E_0=(m_2-m_1)/(m_2+m_1)$，与二元碰撞的标准结果一致，
可作为公式与符号约定的自检。由式可知：测得 $E_1/E_0$ 即可反推 $m_2$（识别元素），
而峰高比可用于推断表面原子的排列与覆盖——这正是 ISS 兼具成分与结构信息的原因。
\end{keycond}

\begin{pitfall}
\textbf{（1）ISS 只看到最外层。}这既是优点也是局限：它无法给出亚表面层的成分，
要获得深度信息必须配合离子溅射做深度剖析，而溅射本身会破坏表面。
\textbf{（2）离子轰击会造成表面损伤与择优溅射。}轻元素优先被溅射，
长时间测量会使表面组成偏离真实值。
\textbf{（3）定量比 XPS/AES 更困难。}散射截面、中和概率与表面几何都影响峰高，
一般只做半定量或相对比较。
\end{pitfall}

\subsection{四种表面方法的比较与选择}

把本章四类方法放在一起看，选择的依据不是「哪个更好」，而是
\textbf{你要的是内层化学态、价带结构、微区成分，还是最外层组成}（表~\ref{tab:surface-compare}）。

\begin{table}[tbp]
\centering
\caption{四类表面分析方法的比较。信息深度一栏决定了「表面」在此处的实际含义\src{据 S2 §9.3--§9.5，书内 p.662--693 / PDF p.153--184 归纳。}}
\label{tab:surface-compare}
{\srcsize
\begin{tabularx}{\twocolwidth}{@{}p{16mm}p{26mm}p{20mm}p{16mm}p{30mm}X@{}}
\toprule
方法 & 探测信号 & 激发源 & 信息深度 & 主要信息 & 主要局限 \\
\midrule
XPS & 光电子（内层） & X 射线 & 约 3--10 nm & 元素组成、化学态、半定量 & 空间分辨较低；荷电与污染影响大；分峰非唯一 \\
UPS & 光电子（价带） & 真空紫外 & 约 1--3 nm（约 10 原子层） & 功函数、价带结构、吸附分子轨道 & 需极洁净表面；绝缘粉末不适用；谱解释难 \\
AES & 俄歇电子 & 电子束 & 约 0.5--2 nm（2--10 原子层） & 轻元素与微区成分、深度剖析 & 重元素谱线重叠；仅导电/半导电样品；束损伤 \\
ISS & 散射离子 & 惰性气体离子 & 最外层原子 & 表面最外层组成与排列 & 仅最外层；离子损伤与择优溅射；定量难 \\
\bottomrule
\end{tabularx}}
\end{table}

\begin{dataflow}
选择顺序：（1）要\textbf{化学态}、且样品可承受 X 射线 → XPS；（2）要\textbf{功函数或
价带}（如金属-载体相互作用、半导体界面） → UPS，且必须先确认表面洁净；
（3）要\textbf{微区}或\textbf{深度剖面}（器件、薄膜、失效分析） → AES；
（4）只关心\textbf{最外层由什么原子终止} → ISS。若三者都需要，
现代多功能谱仪通常可在同一超高真空系统中完成，避免样品在转移中被污染——
\textbf{这是联合使用这三类方法最重要的实际理由}。
\end{dataflow}

\begin{crosscheck}
本章的 XPS 数据处理与常见问题见分片 b；表面分析在催化中的具体应用与准原位技术、
近常压 XPS 见分片 d。若关心的是\textbf{表面吸附物种与反应中间物}而非表面元素组成，
应转向原位振动光谱（\chapref{ch:06}、\chapref{ch:07}）与化学吸附/程序升温技术
（\chapref{ch:05}）——表面分析给出的是「有什么原子」，不是「在怎样反应」。
\end{crosscheck}

\srcnote{本分片来源：S2（下册）§9.3--§9.5，书内 p.662--693 / PDF p.153--184。
式~\eqref{eq:ups-gas}、\eqref{eq:ups-solid} 对应原书式(9-11)、(9-12)，
式~\eqref{eq:aes-quant} 对应式(9-15)，式~\eqref{eq:iss-ch09} 对应式(9-16)；
四式均由「OCR 正文的变量定义 + 学科标准形式」重建，
其中式~\eqref{eq:iss-ch09} 已用 $\theta=90^{\circ}$ 的极限结果做过自检。
原书表格（惰性气体光子能量表 9-23、He 辐射线型表 9-24、AES 灵敏度因子图 9-136、
Cr--Cu 俄歇跃迁谱图 9-137 等）未逐表转录，
需要的读者请按上述页码回查原书。}
